% Template for Affective Computing and Intelligent Interaction (ACII)
%
% Modified 2022-03-15 : Desmond Ong (desmond.c.ong@gmail.com)     
%.      -- update for ACII2022, added section for Ethical Impact Statement
%

\documentclass[conference]{IEEEtran}
\usepackage{upquote}

\IEEEoverridecommandlockouts
% The preceding line is only needed to identify funding in the first footnote. If that is unneeded, please comment it out.
\usepackage{cite}
\usepackage{amsmath,amssymb,amsfonts}
\usepackage{algorithmic}
\usepackage{graphicx}
\usepackage{textcomp}
\usepackage{xcolor}
\def\BibTeX{{\rm B\kern-.05em{\sc i\kern-.025em b}\kern-.08em
    T\kern-.1667em\lower.7ex\hbox{E}\kern-.125emX}}
\begin{document}

\title{A parametric facial expression to represent emotions}
\author{\IEEEauthorblockN{Stefan Dirnstorfer}
\IEEEauthorblockA{\textit{Thetaris GmbH} \\
Munich Germany \\
stefan@thetaris.com}
}
\maketitle

\newcommand{\paramoji}[1]{\hspace*{-1ex}\includegraphics[width=2cm]{gen/paramoji-#1}\hspace*{-1ex}}
\newcommand{\paramojib}[1]{\raisebox{-0.6ex}{\includegraphics[height=2.8ex]{gen/paramojib-#1}}}
\newcommand{\face}[1]{\includegraphics[height=2cm]{gen/#1}}
\newcommand{\faces}[1]{\includegraphics[height=1.6cm]{gen/#1}}

\begin{abstract}
    Two competing approaches have emerged to represent a human emotion in digital form.
    The dimensional approach uses continuous parameters to measure an emotion along independent dimensions,
    such as valence, arousal, and dominance. This is mathematically elegant, but lacks consensus on how
    to represent more complex emotions. The categorical approach uses a discrete number of basic emotions.
    It can capture any emotion that can be described. However, their categorical nature forbids the computation
    of averages, dinstances and other statistical measures.
    To bridge the gap between these two approaches this paper suggests a five dimensional emotional model that
    can be visualized with an emoji style smiley face that can easily be interpreted by humans.
    The resulting parametric facial expressions can represent most of the commonly used emoji faces,
    as well as the seven universal emotions described by Paul Ekman. The resulting graphical
    depictions are clear, expressive and exist with arbitrary nuances and mixtures.
\end{abstract}

\begin{IEEEkeywords}
facial expressions, parametric model, emotion model, visualization, parametric emoji
\end{IEEEkeywords}

\section{Introduction}
Most humans have agreed to use emoticons or emojis with facial expressions
to represent emotions when communicating over a digital
channel.  With skyrocketing occurance on social media platforms there is now a huge amount
of human generated content that is annotated with emotional states.
This could have provided a unique opportunity for humans and computers to exchange emotional
information, as envisioned by Picard\cite{Picard97}. However, this is not happening.
There is currently no known application
where a computer interacts with a human through the meaningful help of emojis.
What makes their emotional content so intuitive for humans but yet so hard to process for computers?

The Unicode Consortium standardized a large body of emojis for computer mediated human-to-human communication.
The standard defines rough gidelines for many emotionally expressive symbols, often in the form
of facial expressions. Unfortunately,
this level of standardization is not sufficient to converge on a globally shared consensus on their semantics.
While the standard may define an emoji as eg. a ``grinning face'', there were counted
at least 17 different versions of how this face may be displayed on various software platforms \cite{emoji_interpretation}.
The consequence is
a severe misunderstanding on even basic attributes of sentiment from positive to negative.
In an anatomy based approach the facial expressions displayed in emojis were studied \cite{emoji_facs}.
An adapted version of the facial action coding system was used to infer muscle activity and the causal emotions,
but exact interpretation varied widely based
on how they were rendered on various platforms and software versions.
The bottom line is that standardized emojis are
not suitable to encode nuanced and precisely defined emotions.

Outside the body of emojis other sets of facial expressions have been graphically
designed to capture fine grained scales of emotions.
Research has found them useful for the communication of qualias, such as pain and taste.
Two effective scales of faces were designed to capture levels of perceived pain within children and adults. \cite{bakerwong, dvprs}.
It was found that visual scales of facial expressions produced more stable and consistent results in self-reported pain and greatly improved pain treatment.
A wider range of emotional experiences can be expressed if additional dimensions are allowed.
The circumplex model of affect is a widely used 2-dimensional model \cite{circumplex} that can be visualized as facial expressions.
An emoji grid with emojis on these to axes was found useful for reporting food elicited emotions\cite{emojigrid}.
Despite the additional depth of two dimensions the grid still lacks space for many emotions.
Additionally, the grid was only defined for emojis on the boundary of the domain without an obvious way to interpolate their visual appearance.
In short, these custom sets are only consistenty defined on very limited domains.
Their expressive range does not nearly cover the variety of Unicode faces.

For a greater range of facial expressions there is a visualization technique that can directly generate comic style
faces from high dimensional data.
Chernoff faces provide unique facial expressions for 18 parameters covering the
sizes of eyes, mouth and ears as well as angles and distances between them\cite{chernoff}.
Unfortunately, this does not include any of the subtle activations that are necessary for the display of deep emotions.
Few other methdods have been tried to sythetically generate realistic facial expressions from
data \cite{parametric_face_drawings, parametric_faces, parametric_face_facs, FACSGen}.
However, without the exaggerated comic style representation the emotional content is difficult to decode.
Therefore, Chernoff faces come closest at being a predecessor, or next best solution, to the one I present in this paper.
Both approaches visualize high dimensional data in a human face. Neither has yet produced a success story in any
practical application. In the following sections I will do my best to make that change.

\section{Five dimensions of the facial expression}
The idea of expressing a human emotion along measurable dimensions is said to date back to Wundt \cite{wundt}, who derived
his ideas mainly through self-inspection.
In a more accessible summary of his work his three dimensions of feeling were translated as pleasantness, excitement and tension\cite{tridimensional}.
The first two of these dimensions have become the universal backbone of emotion classification,
now known as "valence" and "arousal" \cite{circumplex}.
A much more intense debate formed around the third dimension. This paper follows the nomenclature of "dominance"\cite{dominance_mehrabian}.
The fourth dimension is ``contempt'', derived from a discovery that this feeling of superiority is not a distinct emotion, but combines differently with positive and negative feelings. \cite{contempt_ekman}.
As a final dimension I will use ``control', which follows a simplified model for conscious modulation.
Hence, the interpretation of the last dimensions may partly be subject to cultural norms.

It is easy to understand all these dimensions in isolation, but in combination it becomes quickly unclear how to interpret them.
To solve this problem I want to present a visualization in the form of a parametric emoji (paramoji) that provides a
unique and valid facial expression for each combination of parameters.
The face is made to be generic, without any reference to physiological properties, like age, gender, color.
It is the strength of the emotional expression that had the top priority.
With a small exception in contempt, all graphical parameters depended linearly on the emotional dimensions.
Linearity ensures that all paramojis can be smoothly interpolated and animated into each other without jumps or unintended easing.
Although animation is not the topic of this document it is obviously the next frontier of affective computing, because the temporal development of
facial expressions unveils deeper insights that must be parameterized and visualized too.
For all the compromises that had to be made the result should be viewed as a graphical contribution rather than
a psychologial one.

\subsection{Valence and arousal}

The circumplex model of affect is a quantitative framework
to describe all emotional states with two numerical parameters: ``valence'' and ``arousal''\cite{circumplex}.
Valence describes the emotion along the axis of pleasure vs misery and approval vs disapproval. Arousal describes the axis
from aroused to sleepy, active to inactive or agitated to calm.
The claim that all emotions fall into a plane of only two dimensions might appear as quite a stretch.
However, this simplicity makes the model quite attractive and it became the basis of almost every emotional model.
Recent neurological evidence provides further assurance that these dimensions actually exist and correlate to
measurable brain activities. \cite{circumplex2}.

Figure \ref{fig:circumplex} shows the chosen visualization for the two dimensions.
The graphics is defined through the control points of Bezier lines with some rotation of the eyes' outlines.
All coordinates and angles are mapped linearly from the two input dimensions.
In this (still incomplete) model arousal controls eye and mouth sizes and valence controls curvature.

\begin{figure}[htbp!]
    \begin{center}
        \begin{tabular}{cc|ccc}
            && \multicolumn{3}{c}{valence} \\[0.2em]
            && 0.0 & 0.5 & 1.0 \\ \hline
            &&&&\\[-0.8em]
            & \rotatebox{90}{\hspace{1.5em} 1.0} &
            \paramoji{0,100,100,50,0} &
            \paramoji{50,100,100,50,0} &
            \paramoji{100,100,100,50,0} \\
            \rotatebox{90}{\hspace{0.8em} arousal}
            & \rotatebox{90}{\hspace{1.5em} 0.5} &
            \paramoji{0,50,50,50,0} &
            \paramoji{50,50,50,50,0} &
            \paramoji{100,50,50,50,0} \\
            & \rotatebox{90}{\hspace{1.5em} 0.0} &
            \paramoji{0,0,0,50,0} &
            \paramoji{50,0,0,50,0} &
            \paramoji{100,0,0,50,0}
        \end{tabular}
    \end{center}
    \caption{Facial expressions created with three different values for valence ($v$) and arousal ($a$).
    All displacements of graphical control points depend linearly on these two dimensions and can be arbitrarily interpolated. }
    \label{fig:circumplex}
\end{figure}

There is an obvious zero point for arousal when eyes and mouth are totally shut.
On the other end of the spectrum there is no physical limit on how wide a comic eye can open.
Therefore, arousal values can only be compared against each other without an objective center.
For valence the natural center is a flat mouth and flat eyes.
Eyes and mouth bend upwards for pleasure.
This is consistent with a respective muscle movements and was already part of early emoticons as \texttt{\^\_\^} and \texttt{:)}.
Eyes bend downwards for misery. This is derived from the foreshortened view when the head is lowered.
A combination of pleasure with a lowered head, ie a smile with eyes bent downwards, is prevented by the linearity
constraint\footnote{A Unicode emoji for such an expression exists with code 1f60c.}.

Looking at facial expressions one can find the predicted emotions on the edges of the domain:
alarmed \paramojib{50,100,100,50,0},
sleepy \paramojib{50,0,0,50,0},
happy \paramojib{100,50,50,50,0} and
miserable \paramojib{0,50,50,50,0}.
The quadrants also comply with the circumplex model, showing
contentment \paramojib{100,0,0,50,0},
excitement \paramojib{100,100,100,50,0},
distress \paramojib{0,100,100,50,0}
and depression\paramojib{0,0,0,50,0}.
The original publication postulates many more emotions that are supposed to sit tightly packed on a planar space.
However, fear vs anger vs tension, bordom vs sadness: these are fundamentally different emotions that can hardly be seen in close proximity to each other.
More dimensions are needed to separate them.

\subsection{The third dimension: Dominance}

A first candidate for a working three-dimensional model was the evaluation-potency-activity model \cite{potency_osgood}.
It was developed to find the main axes of meaning in nouns by mapping them to adjectives through a rigorous data driven approach.
Factor analysis revealed that two of the dimensions with most explanatory power were equivalent to valence and arousal.
The third of the three most important dimensions of meaning is described by the adjective scales of strong vs weak,
hard vs soft, brave vs cowardly, tenacious vs yielding.
Unfortunately, the author did not find a good name for it.
The pleasure-arousal-dominance model by \cite{dominance_mehrabian} fixed that problem. Its dimensions are
introduced as ``analogues'', because the results of the factor analysis replicated the dimensions with slightly changed names.
In this model anger, power and boldness dominate over anxiety, infatuation, fear and loneliness.

Irrespective of naming and definition details, the third dimension - dominance - shows anger and strength \paramojib{50,50,50,100,0}
on one end of the scale against fear and weakness \paramojib{50,50,50,0,0} on the other end.
The visual effects are lowered inner eyebrows, inwards rotated eyes, more prominent display of teeth and a sharper mouth contour.
For low dominance the visual effects are exactly mirrored, because all displacements are, again, linearly added.
Figure \ref{fig:dominance} shows four of the previous paramojis, which were assumed to have a neutral dominance of 0.5 in the center.
Left and right the paramojis show lowest (0.0) and highest (1.0) values of dominance.
Obviously, the scale is continuous and all intermediate values could be interpolated.

\begin{figure}[htbp!]
    \begin{center}
        \begin{tabular}{l|lccc}
            \rotatebox{90}{\hspace{1ex}$d = 1.0$} &&
            \paramoji{50,50,50,100,0} &
            \paramoji{0,0,0,100,0} &
            \paramoji{0,100,100,100,0}
            \\
            \rotatebox{90}{\hspace{1ex}$d = 0.0$} &&
            \paramoji{50,50,50,0,0} &
            \paramoji{0,0,0,0,0} &
            \paramoji{0,100,100,0,0}
            \\[1ex]
            \hline
            & $v$ & 0.5 & 0.0 & 0.0 \\
            & $a$ & 0.5 & 0.0 & 1.0
        \end{tabular}
    \end{center}
    \caption{Three facial expressions plotted with low (0.0) and strong (1.0) dominance.
    Previous faces above are assumed to be at medium dominance (0.5).
    All expressions can be smoothly interpolated in the interior of the domain where more natural looks are found.}
    \label{fig:dominance}
\end{figure}

The three parameters, valence, arousal and dominance, range from 0 to 1 and span the unit cube domain.
Four corners on the low valence side correspond to misery and disapproval. They already resemble some
preliminary versions of the basic emotions:
fear \paramojib{0,100,100,0,0},
anger \paramojib{0,100,100,100,0},
sadness \paramojib{0,0,0,0,0} and
disgust \paramojib{0,0,0,100,0}.
For positive valence we find pleasure and approval. Two universal emotions sit on the positive end of the cube:
happiness \paramojib{100,50,50,50,0} and
surprise \paramojib{100,100,100,50,0}.
The expressions in the corners of the positive side are probably not universal. I'll call them
simper \paramojib{100,0,0,0,0},
embarrassed \paramojib{100,100,100,0,0},
anticipating \paramojib{100,100,100,100,0}
and a forced smile \paramojib{100,0,0,100,0}.

\subsection{Contempt}

The final missing universal emotion is contempt \paramojib{50,50,0,50,100}.
It is described by the adjectives
scornful, haughty, smug, vain, or disdainful. \cite{contempt_ekman}.
Comtempt is a feeling of moral superiority that varies in expression based on the amount of social support
one gets. Contempt against inferiors, the feeling of power can be a pleasurable experience.
Comtempt against a superior can feel misarable. Because the expression depends on other dimensions it is a dimension
itself, rather than a specific spot in parameters space.

The only universally confirmed expression of contempt is the ``unilateral lip-tightening expression''.
A raised upper lip often shows in conjunction disgust. Tension in the cheek can also press the eyes into
squeezing. A squinting eye is not in itself a signal for contempt,
but a graphical choice to emphasizes its asymmetric nature and,
it often occurs in coincidence.
Figure \ref{fig:contempt} shows random emotions with maximized contempt levels. At least graphically,
contempt mixes perfectly with the previous dimensions. The left eye and the left side of the mouth have a bilinear
graphical dependence. Respective coordinates depend on $c$, but also on pairwise products with the other dimensions.
All combinations of parameters produce comical, but valid facial expressions.

\begin{figure}[htbp!]
    \begin{center}
        \begin{tabular}{l|ccc}
            \rotatebox{90}{\hspace{1ex}$c = 1.0$} &
            \paramoji{86,19,19,68,100} &
            \paramoji{25,48,48,67,100} &
            \paramoji{44,36,36,27,100} \\
            $v$ & 0.86 & 0.25 & 0.44 \\
            $a$ & 0.19 & 0.48 & 0.36 \\
            $d$ & 0.68 & 0.67 & 0.27 \\ \hline
            \rotatebox{90}{\hspace{1ex}$c = 1.0$} &
            \paramoji{81,68,68,33,100} &
            \paramoji{10,14,14,90,100} &
            \paramoji{11,72,72,41,100} \\
            $v$ & 0.81 & 0.10 & 0.11 \\
            $a$ & 0.68 & 0.14 & 0.72 \\
            $d$ & 0.33 & 0.90 & 0.41 \\
        \end{tabular}
    \end{center}
    \caption{Random paramojis with maximum levels of contempt ($c=1$). Graphically, contempt
    combines well with all previous emotions.}
    \label{fig:contempt}
\end{figure}

\subsection{Control}

The final dimension of this model describes the modulation of a facial expression that is partly put under conscious control.
It ranges from controlled, suppressed \paramojib{50,100,0,50,0}
to uncontrolled, expressed \paramojib{50,0,100,50,0}.
The visualization follows an hypothesis raised by Ekman and Friesen\cite{Ekman75}. Their studies suggested that ``when a person is
controlling what is shown on his face, more effort is focused on managing what occurs in and around his mouth and
lips than in the area of the eyes/lids or brows/forehead.''  Although there might be exceptions to that rule,
it is certainly good enough as a first approximation to the subject matter.

The major and easiest to control aspect of an emotion is its intensity. In visual terms, this is arousal ($a$). All
the graphical impact caused by $a$ can be separated into arousal shown in the eye region ($a_1$) and arousal
shown in the mouth region ($a_2$). The amount of excerted control ($o$) can be computed as the difference between
these two. For a consciously deintensified expression $o$ is large. For an intensified expression $o$ is low. The dimensions
$a_1$ and $a_2$ are mathemtically preferred, because they are defined on a cartesian domain. The dimensions $a$ and $o$
are closer to a possible interpretation, but do not produce valid visualizations for all combinations of values.

Equation \eqref{eqn:control} defines how these two dimension pairs can be converted into each other.
The complexity of these mathematical equations is commensurate with the psychological depth of the model.
In the end its the visualization that counts, not the numbers.

\begin{equation}
    \label{eqn:control}
    \begin{array}{rcl}
        a_1 &=& a + o/2 \\
        a_2 &=& a - o/2 \\
        o &=& a_1 - a_2
    \end{array}
\end{equation}

Figure \ref{fig:control} shows the visual effect of the control dimension on some random random faces.
It is interesting that even in the extremely uncontrolled versions one can still differentiate between
positive \paramojib{100,0,100,50,0}
and negative \paramojib{0,0,100,50,0} valence.
This is in contrast to real faces, where extreme pain and extreme joy cannot be distinguished from
the facial expression without context information or body language \cite{painjoy}.
Emoji style visualization are not just exaggerated versions of existing facial expressions.
They can also encode emotional information that go beyond what real faces show.

\begin{figure}[htbp!]
    \begin{center}
        \begin{tabular}{cc|cc}
            \multicolumn{2}{c}{$o = -1.0$} &
            \multicolumn{2}{c}{$o = +1.0$} \\
            \hline
            &&&\\[-0.8em]
            \paramoji{52,0,100,52,0} &
            \paramoji{92,0,100,53,0} &
            \paramoji{9,100,0,52,0} &
            \paramoji{99,100,0,12,92} \\
%   $v$ & 0.52 & 0.92 & 0.09 & 0.99 \\
%   $a_1$ & 0 & 0 & 1 & 1 \\
%   $a_2$ & 1 & 1 & 0 & 0 \\
%   $d$ & 0.52 & 0.53 & 0.52 & 0.12 \\
%   $c$ & 0 & 0 & 0 & 0.92 \\
            \paramoji{27,0,100,69,78} &
            \paramoji{34,0,100,0,0} &
            \paramoji{11,100,0,80,70} &
            \paramoji{65,100,0,70,0} \\
%   $v$ & 0.27 & 0.34 & 0.11 & 0.75 \\
%   $a_1$ & 0 & 0 & 1 & 1 \\
%   $a_2$ & 1 & 1 & 0 & 0 \\
%   $d$ & 0.69 & 0 & 0.80 & 0.79 \\
%   $c$ & 0.78 & 0 & 0.70 & 0 \\
        \end{tabular}
    \end{center}
    \caption{Random emotions with minimal and maximal values of control ($o$). Generally speaking, the mouths is under stronger
    conscious control than the eyes. In a possible overuse of mathematical accuracy,
        control is defined as the difference between respective intensities.}
    \label{fig:control}
\end{figure}

\section{Universal Emotions}

Figure \ref{tab:affnet} visualizes possible parameters for eight basic emotions.
There are seven facial expressions that deviate from the normal face, but only five
dimensions to define the graphical displacements. Hence, there is some compression
of features.
The following is a quick check of the correctness of the visible facial activations against
their original definition provided by Ekman \cite{Ekman72, contempt_ekman}:

\begin{figure}[htbp!]
    \begin{center}
        \begin{tabular}{l|cccc}
            & \hspace*{-1ex} \paramoji{60,50,10,50,0} \hspace*{-2ex}
            & \hspace*{-2ex} \paramoji{100,20,100,50,0} \hspace*{-2ex}
            & \hspace*{-2ex} \paramoji{0,30,0,10,0} \hspace*{-2ex}
            & \hspace*{-2ex} \paramoji{70,90,70,20,0} \hspace*{-1ex}
            \\
            & Normal & Happy & Sad & Surprise \\[1ex]
            $v$ & 0.6 & 1.0 & 0.0 & 0.7 \\
            $a_1$ & 0.5 & 0.2 & 0.3 & 0.9 \\
            $a_2$ & 0.1 & 1.0 & 0.0 & 0.7 \\
            $d$ & 0.5 & 0.5 & 0.1 & 0.2 \\
            $c$ & 0.0 & 0.0 & 0.0 & 0.0 \\ \hline
            & \hspace*{-1ex} \paramoji{40,50,30,20,0} \hspace*{-2ex}
            & \hspace*{-2ex} \paramoji{0,10,60,40,70} \hspace*{-2ex}
            & \hspace*{-2ex} \paramoji{20,60,0,100,0} \hspace*{-2ex}
            & \hspace*{-2ex} \paramoji{60,50,0,70,100} \hspace*{-1ex}
            \\
            & Fear & Disgust & Anger & Contempt \\[1ex]
            $v$ & 0.4 & 0.0 & 0.2 & 0.6 \\
            $a_1$ & 0.5 & 0.1 & 0.6 & 0.5 \\
            $a_2$ & 0.3 & 0.6 & 0.0 & 0.0 \\
            $d$ & 0.2 & 0.4 & 1.0 & 0.7 \\
            $c$ & 0.0 & 0.7 & 0.0 & 1.0
        \end{tabular}
    \end{center}
    \caption{
        Possible parameterizations of the 7 universal emotions.
        All five parameters can be chosen from a continuous domain.
        Many blends and variations could be parameterized and displayed.
        This is just one of many possible representations of the categories.
    }
    \label{tab:affnet}
\end{figure}

\begin{enumerate}
    \item {\bf Happiness:}
    Eyes relaxed and narrowed; lower lid pushed up; outer corners of lips raised; opening mouth and teeth
    appearing; lips are pushed back. The only missing feature are the crow feet at the out edge of the eyes.
    \item {\bf Sadness:}
    Brows drawn together with inner corners raised and outer corners lowered; bulge or muscular
    contraction above center of brow area; drooping upper lid; lax lower lid; lid pulled up at inner corner,
    down at outer corner; eyes look downwards; outer lips pulled down.
    \item {\bf Surprise:}
    Raised curved eye brows; wide open eyes showing sklera above and below the iris; skin above eyelid streched,
    dropped open mouth; lips are parted.
    The weakest feature are the relaxed outer lips that show more tension than optimal.
    \item {\bf Fear:}
    Raised, drawn together eye brows; short horizontal eye brow wrinkles; eyes opened,
    hard stare; the mouth is opened, outer lips are drawn back; sklera on top of the pupil and none below.
    \item {\bf Disgust:}
    Brows drawn down; vertical creases in forhead and mouth; lower eyelid pushed up, but not tense;
    deep nasolabial fold; upper lip raised; lower lip pushed ``outwards'', slightly opening the mouth.
    The weakness is that brows are drawn together too much and there is no wrinkle on the bridge of the nose.
    \item {\bf Anger:}
    Brows pulled down; strong vertical eye brow wrinkles; upper lids appear lowered; lips pressed together
    or showing of teeth. This picture shows the version without teeth. With larger values for $a_2$ teeth would
    appear.
    \item {\bf Contempt:}
    The unilateral lip-tightening is the only signal that is universally confirmed. The choice for valence and arousal
    are actually arbitrary.
\end{enumerate}

\section{A User Interface}

Figure \ref{fig:encoder2} shows a user interface for exploring the parameter space.
Every of the five dimensions has its own slider that ranges from 0 to 100\%. The
influence of each parameter set is unique and can easily be recognized by the resulting shape:
valence controls curvatures; arousal controls the size of eyes and mouth; dominance controls
the slope of the brows; contempt controls horizontal asymmetry; the control dimension leads to
different openings of eyes and mouth. An alternative model uses arousal 1 and arousal 2 instead of
arousal and control to choose the size of eyes and mouth directly.

\begin{figure}[ht!]
    \begin{center}
        \includegraphics[height=5cm]{gen/encode2.pdf}
        \vspace*{-3ex}
    \end{center}
    \caption{
        A possible user interface with sliders for each dimension. A demo version of this interface is
        currently deployed on \texttt{https://paramoji.org/editor.html}.
    }
    \label{fig:encoder2}
\end{figure}


\section{Conclusion}
A lot of effort has been spent on decoding the emotional content from facial expressions.
The suggested approach goes the opposite direction. It creates a facial expression from
emotional dimensions. Unlike ubiquitous emojis the suggested paramojis do not visualize
discrete emotional categories. Instead, they create an emoji style visual representation
from five continuously scalable dimensions. Unlike other dimensional emotional models
the visualization provides a clear and convincing position for the universal emotions
in parameter space.

I believe that the affective computing community would greatly benefit from using the presented
visualizable emotional model. In particular, the application of affect classification from photos of faces
and humans in action such as AffectNet\cite{affectnet} and EMOTIC\cite{emotic_2019} would benefit in
multiple ways. First, a unified framework
for continuous and categorical dimensions would avoid value incongruencies. Second, the
exaggerated visualization makes it much easier to annotate facial expressions and spot potential mistakes.
Third, additional information about
the structure within categories and the representation of border cases would be possible.
Fourth, the use of continuous parameters for all emotions
would enable common statistical post-processing, eg. noise-reduction, filtering etc.
Fifth, affective computing would be much more fun if smileys replaced boring numbers.


\section*{Ethical Impact Statement}

The suggested visualization technique would make it much easier to monitor a detected emotional state.
This could potentially be abused to monitor individual people as well as groups of people.
Another danger is that this simplistic model of an emotion will lead to a reductionist view of emotional needs,
especially when software is developed that detects and aggregates these visualizable emotional dimensions.

\section*{Code}
The source code of the parametric character (paramoji) might seem confusing at first.
However, it is extremely simple,once broken down.
A sparse matrix contains all the parameters. It has a row for each parameter in the drawing instruction and a column
for each entry in the vector $(1, 2v-1, a_1, a_2, 2d-1, c, ca_1,cv, ca_2, cd)$.  The product of the two is inserted
for the question marks in the template. The output format is in SVG, a standardized description language for
vector graphics. Both, the JavaScript code and the SVG instructions, can be interpreted by any contemporary
web browser.

The easiest way to run the code at home is as follows: 1) Open the developer console in your web browser.
2) paste the source code into the console and press enter. 3) Type: \texttt{document.write(paramoji(1,1,0,0,0))}
and press enter again. The facial expression should appear on your screen.

{\small
    \begin{verbatim}
function paramoji(v, a1, a2, d, c) {
 const data=[[51,-2,-10],[,22,,,-26],[,,-11],[-
  8,,-8],[,,-4],[-8,,-8],[,,2],[-8,,-3],[,,9],[
  -4,,1],[,-6,13],[5,,12],[,-9,11],[11,,8],[5,,
  12],[,-9,-11],[-4,,5],[,-6,-13],[,,-11],[,,,,
  -1],[,-2.5,,,-1.5],[.5,,5],[5,1,-4,,8],[2,-9,
  -17,,5],[,-22,,,26,-10,,20],[,,-11,,,,10],[8,
  ,8,,,,-6],[,,-4,,,,4],[8,,8,,,,-6],[,,2,,,,-2
  ],[8,,3,,,,-2],[,,9,,,,-9],[4,,-1],[,-6,13,,,
  -3,-12,6],[-5,,-12],[,-9,11,,,-5,-8,10],[-11,
  ,-8],[-5,,-12],[,-9,-11,,,-5,9,10],[4,,-5],[,
  -6,-13,,,-3,12,6],[,,-11,,,,10],[,,,,2],[2,1,
  -10,,-6,3],[2,-9,-17,,4,20],[,,,,,.4],[70,-2,
  -2.5,-7.5,-1,,,,1],[-3,,.25,.75],[-14,-2,1.5,
  4.5,1],[25,-4,,-6,2,-7],[77,-10,,-2,,-1],[29,
  -2,,-3,2,-9],[76,-1,,-8,-1,-17,-1,1,1],[37,-5
  ,,,2,4],[68,-1,-2,-7,-2,-6,-4,5,,2],[75,4,,6,
  -2],[77,-10,,-2],[71,2,,3,-2],[76,-1,,-8,-1,,
  ,1],[65,4,,-1,-2,-2],[68,-1,-2,-7,-2,-3,-1,3]
  ,[80,5,,6,-6,,,,-2],[66,2,,-5,-1,,,,-2],[14,1
  ,,9,-1],[,-12,,-2,,-3,,6,-2],[10,1,,11,-2],[,
  -2,,11,,,,,-9],[6,,,6,-2],[,1,,12,,,,,-5],[,1
  ,,12,,,,,-4],[-10,-1,,-11,2],[,-2,,11,,,,,-2]
  ,[-14,-1,,-9,1],[,-12,,-2,,-3,,6,2],[-10,-1,,
  -11,2],[,-2,,-11,,,,,-5],[-6,,,-6,2],[,1,,-12
  ,,,,,-3],[,1,,-12,,,,,-1],[10,1,,11,-2],[,-2,
  ,-11,,,,,-2],[14,1,,9,-1],[,-12,,-2,,-3,,6,-2
  ],[,,,,,.4,,.15,-.25],[50,,,,,-12,,,6],[81,2,
  ,,-2,-2,,,1],[59,2,1,,-6],[40,-8,-19,,12],[71
  ,2,3,,-5],[40,-11,-21,,3],[79,1,6,,-4],[40,-3
  ,-18,,-3],[2,-.5,1,,.5],[62,,-4,,-2],[33,-5,-
  16,,7],[56,2,,,-3],[39,-7,-22,,9],[56,1,,,-1]
  ,[41,-8,-20,,10],[,,,,,1,3,2]]
 const template='<svg width="100%" height="10'+
  '0%" viewBox="0 0 100 100"><defs><clipPath '+
  'id="c-eye-r"><use href="#eye-r"/></clipPat'+
  'h><clipPath id="c-eye-l"><use href="#eye-l'+
  '"/></clipPath><clipPath id="c-lips"><use h'+
  'ref="#lips"/></clipPath></defs><g transfor'+
  'm="translate(68,?)"><path id="eye-r" trans'+
  'form="rotate(?)" d="M-8,?C?,? ?,? ?,? ?,? '+
  '?,? ?,0 ?,? ?,? -8,?Z" fill="white" stroke'+
  '="black"/><g clip-path="url(#c-eye-r)"><g '+
  'id="lens" transform="translate(?,?)"><circ'+
  'le r="5" fill="#9d4922"/><circle r="?"/><c'+
  'ircle r="1" cx="-2.5" cy="-1.5" fill="whit'+
  'e"/></g><path d="M-15,-20H20V0Q?,? -15,-2Z'+
  '" opacity=".25"/></g><g transform="transla'+
  'te(-36,0)"><path id="eye-l" transform="rot'+
  'ate(?)" d="M8,?C?,? ?,? ?,? ?,? ?,? ?,0 ?,'+
  '? ?,? 8,?Z" fill="white" stroke="black"/><'+
  'g clip-path="url(#c-eye-l)"><use href="#le'+
  'ns" x="?"/><path d="M-15,-20H15V-2Q?,? -20'+
  ',0Z" opacity=".25"/></g></g></g><path tran'+
  'sform="matrix(1,?,0,1,53,?)" d="M-4,0q-2,-'+
  '2 -4,0M-1,0Q3,-2 4,-1T6,? 4,?" fill="none"'+
  ' stroke="black"/><path d="M?,?Q?,? ?,?" fi'+
  'll="none" stroke="black"/><path d="M?,?Q?,'+
  '? ?,?" fill="none" stroke="black"/><g clip'+
  '-path="url(#c-lips)"><rect height="100" wi'+
  'dth="100"/><ellipse cx="50" cy="91" rx="15'+
  '" ry="10" fill="#800f08"/><g transform="tr'+
  'anslate(0,?)"><use href="#tooth" x="-14"/>'+
  '<use href="#tooth" x="-7"/><use href="#too'+
  'th"/><use href="#tooth" x="7"/></g><g tran'+
  'sform="translate(0,?)"><use href="#tooth" '+
  'transform="matrix(1,.14,0,1,-14,-8)"/><use'+
  ' href="#tooth" x="-7"/><rect id="tooth" x='+
  '"50.5" rx="2" ry="1" height="15" fill="whi'+
  'te" width="6" stroke="black" stroke-width='+
  '".5"/><use href="#tooth" transform="matrix'+
  '(1,-.14,0,1,7,7)"/></g></g><path id="lips"'+
  ' d="M?,?C?,? ?,? 0,?S?,? ?,?C?,? ?,? 0,?S?'+
  ',? ?,?Z" transform="matrix(1,?,0,1,?,?)" f'+
  'ill="none" stroke="black"/><g id="brow-r">'+
  '<path d="M?,?Q?,? ?,?" stroke-width="?" st'+
  'roke="black" fill="none"/><path d="M?,?Q?,'+
  '? ?,?" fill="none" stroke="black"/></g><us'+
  'e href="#brow-r" transform="matrix(-1,0,0,'+
  '1,100,?)"/></svg>'
 const dotprod = (X,Y) => X.reduce(
  (a,b,i) => a + b*(Y[i] || 0), 0)
 let index = 0,
  V = [1,2*v-1,a1,a2,2*d-1,c,c*a1,c*v,c*a2,c*d]
 return template.replace(/\?/g,
  ()=> dotprod(V, data[index++]))
 }
    \end{verbatim}
}


\bibliographystyle{plain}
\bibliography{paramoji}

\end{document}
